\section{Tecnologias}

O desenvolvimento do sistema \textbf{Tropical Mix} foi estruturado a partir de um conjunto de tecnologias modernas e amplamente utilizadas no mercado de desenvolvimento web. A escolha das ferramentas considerou critérios como desempenho, escalabilidade, segurança e facilidade de manutenção. A seguir, são descritas as principais tecnologias aplicadas em cada camada da aplicação.

\subsection{Front-end}

A camada de \textit{front-end} do sistema foi desenvolvida utilizando a biblioteca \textbf{React.js}, que possibilita a criação de interfaces de usuário modernas, dinâmicas e responsivas. Essa tecnologia, baseada em componentes reutilizáveis, permite a construção de telas interativas e modulares, facilitando a manutenção e a evolução da aplicação.  

A escolha pelo React deve-se à sua eficiência no gerenciamento do estado da aplicação, à sua capacidade de atualização em tempo real e à ampla comunidade de desenvolvedores, o que assegura suporte contínuo e uma vasta gama de bibliotecas complementares.  

Para a estruturação e estilização das páginas, foram utilizadas as tecnologias \textbf{HTML5} e \textbf{CSS3}, permitindo a criação de layouts responsivos que se adaptam a diferentes dispositivos, como computadores, tablets e smartphones. A comunicação entre o \textit{front-end} e o \textit{back-end} é realizada por meio de requisições HTTP utilizando a biblioteca \textbf{Axios}, que possibilita a troca de dados de forma assíncrona e eficiente.  

Além das funcionalidades de navegação e controle de estoque e vendas, o \textit{front-end} também inclui o módulo de um painel dinâmico, \textbf{Dashboard}, um painel interativo voltado à geração de relatórios e à exibição de indicadores de desempenho. Esse painel consome dados disponibilizados pela API do \textit{back-end} e apresenta informações estratégicas, como volume de vendas, produtos com menor estoque e desempenho por período, em gráficos e tabelas dinâmicas.  

Dessa forma, o \textit{front-end} não apenas oferece uma interface intuitiva e agradável ao usuário, mas também serve como ferramenta de apoio à gestão, centralizando visualmente as informações mais relevantes para a tomada de decisão.

\subsection{Back-end}

O \textit{back-end} foi implementado utilizando a linguagem \textbf{Python}, em conjunto com o microframework \textbf{Flask}, responsável por gerenciar as rotas, processar as regras de negócio e realizar a comunicação com o banco de dados. O Flask foi escolhido por sua leveza, flexibilidade e fácil integração com outras tecnologias, além de oferecer suporte nativo à construção de APIs RESTful.  

Entre as principais funcionalidades desenvolvidas nessa camada, destacam-se:
\begin{itemize}
    \item Gerenciamento de usuários e autenticação, com controle de perfis de acesso;
    \item Registro e controle de vendas e movimentações de estoque;
    \item Processamento de dados para geração de relatórios e indicadores exibidos no painel do \textit{front-end};
    \item Integração com o banco de dados \textbf{MySQL} para armazenamento e consulta das informações.
\end{itemize}

O uso do Flask proporciona um desenvolvimento ágil, modular e seguro, permitindo que novas funcionalidades sejam incorporadas de forma independente, sem comprometer a arquitetura do sistema. A comunicação com o \textit{front-end} ocorre por meio de uma \textbf{API REST}, garantindo a troca estruturada e eficiente de dados no formato \textbf{JSON}.


\subsection{Banco de Dados}

O banco de dados escolhido para o sistema é o \textbf{MySQL}, um dos gerenciadores de banco de dados relacionais mais consolidados no mercado. Sua escolha foi motivada pela confiabilidade, escalabilidade e ampla compatibilidade com aplicações baseadas em Python.  

O modelo de dados foi projetado de forma a garantir integridade referencial e desempenho nas operações de leitura e escrita. Cada transação, seja de venda, compra ou atualização de estoque, é registrada no banco de dados, permitindo o acompanhamento detalhado das atividades da loja.  

Além disso, a estrutura relacional do MySQL possibilita consultas otimizadas para a geração de relatórios, facilitando a análise de dados e o processo de tomada de decisão por parte da gestão da \textbf{Tropical Mix}.



\subsection{Infraestrutura}

TEREMOS NÉ?


\subsection{Ferramentas de apoio}

TERAÁ ????